واصل معدل نمو الناتج المحلي الإجمالي ارتفاعه مسجلا 5.3\% في الربع الأول (يوليو - سبتمبر) من العام المالي 2026/2025، وهو أعلى مستوى له في ثلاث سنوات، مما يعكس استقرار الأوضاع الاقتصادية الكلية. بالإضافة إلى بدء تعافي قطاعات الصناعة التحويلية غير البترولية، والسياحة والاتصالات وتكنولوجيا المعلومات. فيعكس ارتفاع النمو بعد هدوء التوترات الجيوسياسية في البحر الأحمر. ومن جانب الانفاق، تحسن مساهمة كل من الاستثمار الخاص وصافي الصادرات في الناتج خلال الربع الأول من العام المالي 2026/2025 (وزارة التخطيط والتنمية الاقتصادية والتعاون الدولي، 2025).\textsuperscript{8}
